%% ISMB / Bioinformatics (OUP) submission template
%% Title: Systematic benchmarking of functional module partitioning in GO pathway subnetworks
%% Compile with: pdflatex -> bibtex -> pdflatex -> pdflatex
\documentclass[twocolumn,11pt]{article}

\usepackage[utf8]{inputenc}
\usepackage[T1]{fontenc}
\usepackage{times}
\usepackage{graphicx}
\usepackage{amsmath,amssymb}
\usepackage{booktabs}
\usepackage{multirow}
\usepackage[numbers,sort&compress]{natbib}
\usepackage[colorlinks=true,linkcolor=blue,citecolor=blue,urlcolor=blue]{hyperref}
\usepackage[margin=2.0cm]{geometry}
\usepackage{caption}
\captionsetup{font=small,labelfont=bf,skip=4pt}
\usepackage{titlesec}
\titlespacing*{\section}{0pt}{8pt}{4pt}
\titlespacing*{\subsection}{0pt}{6pt}{3pt}

\title{\vspace{-1.2cm}\bfseries Systematic benchmarking of functional module partitioning
in GO pathway subnetworks: multi-resolution community detection versus
hierarchical co-expression clustering}

\author{Mingyu Sun\textsuperscript{1,*}\\[2pt]
\small \textsuperscript{1}Affiliation to be completed\\
\small \textsuperscript{*}To whom correspondence should be addressed.}

\date{}

\begin{document}
\maketitle

\begin{abstract}
\noindent\textbf{Motivation.} Weighted gene co-expression network analysis (WGCNA) and its
variants remain the dominant paradigm for partitioning genes into functional modules.
However, the assumption that hierarchical clustering on topological overlap adequately
represents functional structure within a single Gene Ontology (GO) pathway subnetwork has
not been systematically evaluated. Three issues remain unaddressed: whether dynamic tree
cutting resolves nested functional sub-structures; whether linear distance assumptions
conflate functionally antagonistic genes; and whether reported modularity reflects genuine
biological organisation rather than an artefact of network construction.

\vspace{4pt}
\noindent\textbf{Results.} We assembled a benchmark of 11 publicly available expression
cohorts spanning neoplastic and non-neoplastic disease, seven GO biological-process
pathways and 31 partitioning algorithms from eight methodological families, yielding 70
pathway--cohort pairs and 2170 evaluation records. All algorithms operated on an identical
topological overlap matrix per pair, enforcing a strictly paired comparison. Evaluation
followed a six-layer protocol covering structural compactness, functional coherence
(GO/KEGG/Reactome-referenced, size-unbiased $z$-scores), bootstrap stability,
multi-perturbation robustness, biological validity (curated antagonistic gene pairs) and
downstream clinical prediction. To eliminate circularity, graph representation learning was
reformulated as self-supervised edge prediction, removing all community labels from
training. Synthetic benchmarks with known ground truth supplied absolute accuracy
references. Multi-resolution community detection produced functionally more coherent
modules than hierarchical clustering (median $z = 0.768$ versus $-0.075$; paired Wilcoxon
$q = 5.3\times10^{-7}$; Cliff's $\delta = +0.522$) while simultaneously improving modularity
($0.304$ versus $0.000$; $q = 1.8\times10^{-12}$), and recovered nested functional
sub-structures with markedly higher accuracy in synthetic benchmarks (ARI $1.000$ versus
$0.908$). Downstream clinical prediction also favoured multi-resolution partitions
(AUC $0.989$ versus $0.898$).

\vspace{4pt}
\noindent\textbf{Conclusions.} Two methodological caveats temper the interpretation of
modularity-based evidence: observed modularity did not exceed degree-preserving null models
($z = -5.54$), and self-supervised embeddings only partially reproduced the inferred
communities. We recommend reporting functional coherence against pathway-internal null
expectations combined with downstream predictive utility, rather than modularity or module
counts alone.

\vspace{6pt}
\noindent\textbf{Contact:} [email] \quad
\textbf{Supplementary information:} available at [URL]\\
\noindent\textbf{Keywords:} functional module partitioning; community detection; Leiden
algorithm; WGCNA; graph neural networks; self-supervised learning; benchmarking
\end{abstract}

\vspace{2pt}

%% ============================================================
\section{Introduction}
%% ============================================================

Genes do not act in isolation. Their products participate in overlapping complexes and
pathways, and the analysis of coordinated expression across genes is a cornerstone of
systems biology. A dominant analytical strategy treats co-expression as a network and
partitions it into modules, under the assumption that genes grouped together share
function and that module-level summaries (eigengenes, hubs) are biologically meaningful
\citep{langfelder2008wgcna}.

This paradigm has produced substantial biological insight, and WGCNA has become the
\emph{de facto} standard for module identification in bulk transcriptomic studies.
However, its application to the restricted setting of a single functional pathway --- a
routine step in pathway-centric analyses and in hypothesis generation for drug-target
discovery --- rests on assumptions that have received comparatively little scrutiny:

\begin{enumerate}
\item \textbf{Linearity.} Hierarchical clustering on topological overlap dissimilarity
presumes that functional relatedness is adequately captured by a monotone transformation
of co-expression.
\item \textbf{Global cutting.} Dynamic tree cutting determines module boundaries from a
global dendrogram height, rendering module assignment sensitive to overall network size.
\item \textbf{Single scale.} WGCNA produces a single partition. Whether an identified
module contains coherent sub-structures at finer scales is not addressed by the method.
\item \textbf{Absence of external validation.} There is no built-in mechanism to test
whether the inferred partition reflects genuine structure.
\end{enumerate}

Furthermore, when a pathway is defined by a GO term, it is by construction a \emph{nested}
hierarchy: ``regulation of immune response'' contains ``cellular response to cytokine
stimulus'', each of which contains coherent sub-processes. A method that outputs a single
coarse partition discards this hierarchy by design.

\subsection{Related work}

Multi-resolution community detection, in particular the Leiden algorithm
\citep{traag2019leiden}, generalises modularity optimisation through a resolution
parameter and guarantees well-connected communities. Markov Stability provides an
equivalent formulation in continuous Markov time \citep{delvenne2010stability}.
Alternative approaches include the Constant Potts Model \citep{traag2011cpm}, which
eliminates the resolution limit of modularity; information-theoretic methods
\citep{rosvall2008infomap}; generative models including degree-corrected and hierarchical
stochastic block models; overlapping community detection \citep{ahn2010link,
lancichinetti2011oslom}; and graph representation learning, which embeds nodes and derives
partitions in the embedding space. Within bioinformatics, MEGENA \citep{song2015megena}
explicitly addresses multi-scale co-expression module discovery.

Nevertheless, to our knowledge no study has (i) compared partitioning algorithms on
identical pathway-derived graphs across disease categories with a strictly paired design,
(ii) combined synthetic benchmarks with functional-coherence metrics and downstream
clinical tasks, and (iii) reported null-model results that constrain the interpretation of
modularity.

\subsection{Contributions}

\begin{enumerate}
\item A \textbf{paired benchmark design} in which 31 partitioning algorithms operate on an
identical topological overlap matrix per pathway--dataset pair.
\item A \textbf{six-layer evaluation protocol} separating structural, functional,
stability, robustness, biological and downstream dimensions, and explicitly avoiding
module-count-based criteria.
\item \textbf{Functional coherence with a pathway-internal null}, replacing enrichment
$p$-values, which are biased by module size, with a size-unbiased $z$-score.
\item A \textbf{self-supervised reformulation} of graph representation learning, removing
the label leakage inherent in training an embedding model on the community labels it is
later used to validate.
\item \textbf{Synthetic benchmarks with absolute ground truth}, including a nested
sub-structure benchmark that directly tests whether multi-resolution partitioning recovers
fine-grained functional divisions.
\item \textbf{Two negative results}, reported without suppression: observed modularity
does not exceed degree-preserving null models, and self-supervised embeddings do not fully
reproduce community labels.
\end{enumerate}

%% ============================================================
\section{Materials and methods}
%% ============================================================

\subsection{Datasets}

Eleven publicly available cohorts spanning neoplastic and non-neoplastic disease were
retrieved from the Gene Expression Omnibus (GEO) (Table~\ref{tab:datasets}). Each cohort
was processed to a gene $\times$ sample matrix (gene symbols, uppercase), filtered to
remove low-expression genes by an adaptive rule (a gene was retained if its expression
exceeded the within-cohort 25th-percentile-median threshold in more than 50\% of samples,
with progressively relaxed thresholds until at least 1000 genes remained), and
log$_2$-transformed where the dynamic range exceeded 100. Probe-to-gene mapping used the
mygene.info annotation service. Datasets whose series matrix contained no expression
values were excluded.

Within each cohort, seven core GO biological-process pathways were analysed: apoptotic
process (GO:0006915), inflammatory response (GO:0006954), regulation of immune response
(GO:0050776), response to cytokine (GO:0034097), regulation of I-$\kappa$B
kinase/NF-$\kappa$B signalling (GO:0043122), cellular response to oxidative stress
(GO:0034599) and chemokine-mediated signalling pathway (GO:0070098). A pathway was
analysed only if at least 25 of its genes had measurable expression; pathways larger than
300 genes were reduced to the 300 highest-variance genes to bound computation.

\subsection{Network construction}

For each pathway--cohort pair, a soft-thresholded correlation adjacency was computed as
$a_{ij} = |\mathrm{cor}(x_i,x_j)|^{\beta}$ with $\beta = 6$, and the topological overlap
matrix (TOM) derived as
\begin{equation}
\mathrm{TOM}_{ij} = \frac{l_{ij} + a_{ij}}{\min(k_i,k_j) + 1 - a_{ij}},
\qquad l_{ij} = \sum_u a_{iu}a_{uj},
\end{equation}
where $k_i = \sum_u a_{iu}$. The TOM served as the \emph{single shared input} to every
algorithm, so all differences in output are attributable to the partitioning procedure
rather than to graph construction.

\subsection{Partitioning algorithms}

Thirty-one algorithms spanning eight families were evaluated (Table~\ref{tab:algos}).

\noindent\textbf{Family A (hierarchical co-expression).} Average-linkage hierarchical
clustering with dynamic-style cutting (balance criterion), serving as the reference.

\noindent\textbf{Family B (modularity optimisation).} Leiden across a resolution grid
($\gamma = 0.1, 0.5, 1.0, 2.0, 5.0, 10.0$), Louvain ($\gamma = 0.5, 1.0, 2.0$), Constant
Potts Model ($\gamma = 0.05, 0.10, 0.20$), Multilevel, Fast Greedy
\citep{clauset2004fastgreedy}, Leading Eigenvector and Spin Glass.

\noindent\textbf{Family C (information-theoretic / spectral).} InfoMap, Walktrap, spectral
clustering on the precomputed affinity.

\noindent\textbf{Family D (generative / label propagation).} Label propagation.

\noindent\textbf{Family E (overlapping community detection).} Clique percolation.

\noindent\textbf{Family F (graph representation learning).} Self-supervised graph attention
networks (GAT) and GraphSAGE, with partitions derived by $k$-means in the embedding space
(Section~\ref{sec:selfsup}).

\noindent\textbf{Family G (classical clustering).} $K$-means, Gaussian mixture, Ward
agglomerative clustering, DBSCAN, HDBSCAN, OPTICS, Mean Shift, Affinity Propagation and
non-negative matrix factorisation. Density-based methods operated in a spectral embedding
space obtained by Laplacian eigenmaps of the TOM, to avoid the curse of dimensionality
that affects density estimation on high-dimensional similarity profiles.

\noindent\textbf{Family H (biological prior).} Not applicable in the present design;
pathway membership is the grouping prior and is shared by all algorithms.

All stochastic algorithms used a fixed seed (42); stability was assessed across bootstrap
resamples rather than across seeds alone.

\subsection{Self-supervised graph representation learning}
\label{sec:selfsup}

To avoid circularity --- an embedding model trained on community labels cannot provide
independent evidence about those labels --- the representation component was reformulated
as self-supervised edge prediction. For a graph with node features $X$ (standardised rows of
the TOM) and weighted edge set $E$, a fraction (15\%) of edges was masked; a two-layer
graph attention network with edge-feature conditioning (\texttt{edge\_dim}$=1$) produced
$d = 64$-dimensional node embeddings; and the objective was a binary cross-entropy over
positive (masked) edges and an equal number of uniformly sampled negatives:
\begin{equation}
\mathcal{L} = -\frac{1}{|E^{+}|}\sum_{(i,j)\in E^{+}}\log\sigma(z_i^{\top}z_j)
-\frac{1}{|E^{-}|}\sum_{(i,j)\in E^{-}}\log\!\left(1-\sigma(z_i^{\top}z_j)\right),
\end{equation}
with early stopping (patience 40). Because community labels are never used during training,
the degree to which the Leiden partition is separable in the embedding space constitutes a
genuinely independent test. A control experiment on a synthetic graph with known community
structure confirmed that the implementation recovers ground-truth partitions (ARI $=1.000$).

\subsection{Evaluation protocol}

Six orthogonal layers were computed for each algorithm--pathway--cohort triple.

\noindent\textbf{L1 (structure).} Number of modules, Newman modularity $Q$, module-size
entropy, unassigned-gene rate.

\noindent\textbf{L2 (functional coherence; primary functional endpoint).} For each module,
the pairwise Jaccard similarity of GO/KEGG/Reactome annotations (6807 gene sets with
$\geq 10$ genes) was averaged over annotated gene pairs. The module score was standardised
against 30 random gene sets of identical size drawn from the \emph{same pathway}:
\begin{equation}
z_m = \frac{c_m - \mathbb{E}[c_{\mathrm{null}}]}{\mathrm{sd}[c_{\mathrm{null}}]}.
\end{equation}
This removes both module-size bias and the pathway-level enrichment background. Enrichment
$p$-values were retained only as a secondary, size-biased reference.

\noindent\textbf{L3 (stability).} Bootstrap resampling of 80\% of genes, 4 repeats; mean
pairwise ARI across resamples.

\noindent\textbf{L4 (robustness).} ARI between the baseline partition and partitions
obtained after adding 10\% and 20\% Gaussian noise to the TOM and after removing 10\% of
edges.

\noindent\textbf{L5 (biological validity).} Twenty literature-curated antagonistic gene
pairs (e.g. BAX/BCL2, CASP3/XIAP, NFKB1/NFKBIA); the separation rate is the fraction
assigned to different modules.

\noindent\textbf{L6 (downstream clinical prediction).} Module eigengenes (first principal
component of each module) were used as features to predict clinically annotated phenotypes
by logistic regression under five-fold stratified cross-validation, reporting the area
under the ROC curve (AUC).

\noindent\textbf{Random control.} ARI between the observed partition and a uniformly random
partition with the same number of modules.

\subsection{Synthetic benchmarks with ground truth}

\noindent\textbf{B1 (block model).} Weighted graphs with seven known communities, generated
with a fuzziness parameter $\mu\in\{0.10,0.25,0.40,0.55,0.70\}$ controlling overlap between
within- and between-community edge weights; three seeds per setting.

\noindent\textbf{B2 (nested sub-structure; decisive).} Six modules, each containing two
functional sub-groups (12 ground-truth groups), with strong within-sub-group, weak
between-sub-group and near-zero between-module weights.

\noindent\textbf{B3 (noise gradient).} Block-model graphs ($\mu=0.25$) subjected to
Gaussian weight noise $\sigma\in\{0,0.10,0.20,0.35,0.50\}$.

\subsection{Null models}

\noindent\textbf{Null A (degree-preserving rewiring).} Edges were rewired while preserving
each node's degree ($3|E|$ swaps), then re-weighted by sampling from the observed weight
distribution; Leiden modularity was recomputed on each of 30 realisations.

\noindent\textbf{Null B (expression permutation).} Each gene's expression vector was
independently permuted across samples, the TOM recomputed and modularity evaluated as above.

\subsection{Statistical analysis}

The unit of analysis is a pathway--cohort pair. Because all algorithms received an identical
input graph, comparisons are paired by construction. Paired differences were assessed with
the Wilcoxon signed-rank test, corrected for multiple comparisons by the
Benjamini--Hochberg procedure \citep{benjamini1995fdr} ($q<0.05$), and accompanied by
Cliff's $\delta$ as a non-parametric effect size. Medians are reported for skewed metrics.

\subsection{Software and reproducibility}

Analyses used Python 3.11 with \texttt{igraph}, \texttt{leidenalg}, \texttt{networkx},
\texttt{scikit-learn}, \texttt{PyTorch Geometric}, \texttt{hdbscan} and \texttt{mygene}.
All random seeds were fixed. Scripts, processed matrices and result tables are organised as
a linear pipeline (acquisition $\rightarrow$ pathway extraction $\rightarrow$ graph
construction $\rightarrow$ partitioning $\rightarrow$ embedding $\rightarrow$ evaluation
$\rightarrow$ statistics $\rightarrow$ visualisation $\rightarrow$ reporting).

%% ============================================================
\section{Results}
%% ============================================================

\subsection{Benchmark scale}

Eleven cohorts were retained after quality control (Table~\ref{tab:datasets}): three lung
cancer, two breast cancer, two colorectal cancer and one each of hepatocellular carcinoma,
sepsis, type 2 diabetes, cardiomyopathy and systemic lupus erythematosus. Seven GO
biological-process pathways were analysed per cohort where at least 25 genes had measurable
expression, yielding 70 analysable pathway--cohort pairs (median 177 genes per pathway).
Thirty-one partitioning algorithms produced 2170 evaluation records.

\subsection{Multi-resolution partitioning improves functional coherence at equal or better
modularity}
\label{sec:main}

Paired comparison of Leiden at $\gamma=1.0$ against WGCNA-style hierarchical clustering
across all 70 pathway--cohort pairs showed that multi-resolution partitioning produced
significantly more modules (median 4.5 versus 3.0; Wilcoxon signed-rank,
$q = 2.1\times10^{-11}$, Cliff's $\delta = +0.857$) while simultaneously achieving higher
modularity (median $Q = 0.304$ versus $\mathbf{0.000}$; $q = 1.8\times10^{-12}$,
$\delta = +0.989$). The reference method's median modularity of exactly zero indicates that
hierarchical cutting on pathway-derived graphs frequently fails to yield a partition that
improves on the trivial one-module solution.

Crucially, the same comparison on the size-unbiased functional coherence metric --- computed
against random gene sets drawn from the same pathway --- favoured multi-resolution
partitioning (median $z = 0.768$ versus $-0.075$; $q = 5.3\times10^{-7}$,
$\delta = +0.522$). This establishes that the finer partitions are not merely more numerous
but functionally more coherent than the pathway's own baseline expectation
(Fig.~\ref{fig:metrics}, Fig.~\ref{fig:heatmap}).

The separation rate for 20 literature-curated antagonistic gene pairs was also higher for
multi-resolution partitioning (median $0.500$ versus $0.000$; $q = 9.5\times10^{-5}$,
$\delta = +0.633$, 31 pairs with testable gene pairs). Stability did not differ
significantly (median $0.642$ versus $0.635$; $q = 0.495$).

\subsection{Resolution controls a precision--fragmentation trade-off}

Sweeping the resolution parameter from $\gamma = 0.1$ to $\gamma = 10.0$ revealed a
systematic trade-off (Fig.~\ref{fig:overpartition}). Module count increased steeply
(1.2 $\rightarrow$ 83.5 modules), but modularity collapsed ($0.304$ at $\gamma=1.0$
$\rightarrow$ $0.050$ at $\gamma=10.0$) and functional coherence declined. The intermediate
range $\gamma\in[1.0,2.0]$ occupied the region of maximal functional coherence, whereas
$\gamma\geq 5.0$ produced partitions with modularity approaching zero. Module count alone is
therefore an inadequate criterion, and resolution must be chosen against an external
functional criterion.

\subsection{Synthetic benchmarks confirm recovery of nested structure}

On the block-model benchmark with known ground truth, most algorithms recovered the planted
communities when fuzziness was low, but performance diverged sharply as the between- and
within-community weight distributions overlapped ($\mu = 0.55$--$0.70$). Leiden at
$\gamma=1.0$ and Louvain were the most robust, whereas high-resolution
($\gamma = 5.0$) and low-resolution ($\gamma = 0.1$, CPM $\gamma = 0.05$) settings failed
systematically.

The decisive test was the nested sub-structure benchmark, in which six modules each
contained two functionally distinct sub-groups (12 ground-truth groups). Leiden at
$\gamma = 2.0$, CPM at $\gamma = 0.10$, InfoMap and label propagation recovered the
fine-grained partition exactly (ARI $=1.000$), whereas WGCNA-style hierarchical clustering
reached ARI $= 0.908$ and identified only 11 modules. Low-resolution settings
($\gamma=0.5$, spectral clustering) merged the sub-groups entirely (ARI $= 0.602$,
6 modules). This provides direct, ground-truth-based evidence that multi-resolution
partitioning resolves nested functional structure that single-scale hierarchical cutting
cannot (Fig.~\ref{fig:synthetic}).

\subsection{Self-supervised embeddings provide independent structural support that
favours multi-resolution partitioning}

Because the embedding model was trained exclusively on an edge-prediction objective with no
access to community labels, agreement between the learned representation and a partition
constitutes an independent test of whether that partition corresponds to structure
recoverable from the graph alone. On the apoptotic subnetwork, the Leiden $\gamma=1.0$
partition agreed with the embedding-derived partition substantially more than the
hierarchical reference did (ARI $= 0.298$ versus $0.048$), while Leiden $\gamma=2.0$
occupied an intermediate position (ARI $= 0.162$) and the two Leiden solutions agreed only
weakly with the hierarchical partition (ARI $= 0.015$) (Fig.~\ref{fig:umap}). The same
ordering held across pathway--cohort pairs. Two implications follow. First, the low absolute
agreement confirms that no partition is fully recoverable by a shallow self-supervised
encoder, so community structure should not be over-interpreted. Second, and more
importantly, the \emph{relative} agreement is informative: an embedding that never observed
any community label nonetheless reproduces the multi-resolution partition several-fold
better than the hierarchical one, providing independent structural corroboration for the
main result of Section~\ref{sec:main}. On a synthetic graph with known communities, the same
pipeline recovered the ground-truth partition exactly (ARI $=1.000$).

\subsection{Downstream clinical prediction}

Module eigengenes derived from each partition were used to predict disease status under
five-fold stratified cross-validation (317 records). Multi-resolution partitions achieved
consistently higher AUC than the hierarchical reference (Leiden $\gamma=1.0$: $0.989$;
Leiden $\gamma=2.0$: $0.991$; Louvain $\gamma=2.0$: $0.992$; WGCNA-style: $0.898$)
(Fig.~\ref{fig:downstream}). While absolute AUC values are high for all methods ---
reflecting strong transcriptomic separation between disease and control tissue --- the
relative ordering is consistent with the functional-coherence results. Notably, algorithms
producing very fine partitions (e.g. CPM $\gamma=0.20$, $>140$ modules) did not outperform
moderate-resolution settings, again indicating that excessive fragmentation is not
beneficial.

\subsection{Null models constrain the interpretation of modularity}
\label{sec:null}

Both null models produced modularity values \emph{higher} than those observed on the real
pathway networks (degree-preserving rewiring: mean $z = -5.54$; expression permutation: mean
$z = -5.09$). That is, the observed modularity did not exceed what degree-preserving
randomisation of the same graph attains (Fig.~\ref{fig:null}). This is a substantive
negative result: modularity $Q$ on pathway-derived co-expression graphs should not be used
as evidence that inferred modules reflect organisation beyond what the network's degree
distribution already implies. We therefore base our primary conclusions on functional
coherence and downstream predictive utility, both defined with respect to explicit external
references (annotation databases and clinical phenotypes).

This result also illuminates the comparison in Section~\ref{sec:main}: the hierarchical
reference's median modularity of zero is not merely ``worse'' than Leiden's $0.304$;
rather, modularity itself carries little biological information in this setting.

%% ============================================================
%% FIGURES
%% ============================================================

\begin{figure}[t]
\centering
\includegraphics[width=\columnwidth]{../output/figures_v3/Fig2_metrics.png}
\caption{Multi-metric comparison of the 31 partitioning algorithms. (a) Functional
coherence $z$-score ranking. (b) Stability (bootstrap ARI) versus robustness (ARI under
10\% noise). (c) Separation rate for literature-curated antagonistic gene pairs.}
\label{fig:metrics}
\end{figure}

\begin{figure}[t]
\centering
\includegraphics[width=\columnwidth]{../output/figures_v3/Fig3_heatmap.png}
\caption{Performance profile of each algorithm across six evaluation layers ($z$-score
standardised). The primary functional endpoint (coherence $z$) is shown alongside
structural, stability, robustness and biological metrics.}
\label{fig:heatmap}
\end{figure}

\begin{figure}[t]
\centering
\includegraphics[width=\columnwidth]{../output/figures_v3/Fig1_overpartition.png}
\caption{Fragmentation cost of resolution. (a) Number of modules versus modularity $Q$.
(b) Number of modules versus functional coherence. (c) Continuous sweep of the resolution
parameter, showing simultaneous collapse of modularity and coherence at high resolution.}
\label{fig:overpartition}
\end{figure}

\begin{figure}[t]
\centering
\includegraphics[width=\columnwidth]{../output/figures/Fig5_synthetic.png}
\caption{Synthetic benchmarks with known ground truth. (a) Absolute accuracy (ARI) as a
function of community fuzziness. (b) Recovery of nested sub-groups (the decisive
benchmark). (c) Robustness to additive noise.}
\label{fig:synthetic}
\end{figure}

\begin{figure}[t]
\centering
\includegraphics[width=\columnwidth]{../output/figures/Fig7_downstream.png}
\caption{Downstream clinical prediction. Cross-validated AUC for module-eigengene
classifiers predicting disease status, shown per algorithm and as a function of module
count.}
\label{fig:downstream}
\end{figure}

\begin{figure}[t]
\centering
\includegraphics[width=\columnwidth]{../output/figures/Fig6_null.png}
\caption{Null-model comparison. (a) Observed modularity against degree-preserving rewired
and expression-permuted null models. (b) Distribution of $z$-scores relative to the null
models.}
\label{fig:null}
\end{figure}

\begin{figure}[t]
\centering
\includegraphics[width=\columnwidth]{../output/figures_v3/Fig4_consistency.png}
\caption{Cross-cohort consistency. $\Delta$ functional coherence (Leiden $\gamma=1.0$ minus
WGCNA-style) per cohort, and the distribution of coherence $z$-scores for both methods.}
\label{fig:consistency}
\end{figure}

%% ============================================================
%% TABLES
%% ============================================================

\begin{table*}[t]
\centering
\caption{Discovery cohorts. All datasets are publicly available from GEO.}
\label{tab:datasets}
\small
\begin{tabular}{lllrr}
\toprule
Accession & Disease & Tissue & Samples & Pathway genes (median) \\
\midrule
GSE10072 & Lung adenocarcinoma & Lung & 107 & 177 \\
GSE137342 & Severe sepsis & Whole blood & 57 & 168 \\
GSE15653 & Type 2 diabetes / obesity & Liver & 18 & 155 \\
GSE17536 & Colorectal cancer & Colon & 177 & 164 \\
GSE19188 & Non-small cell lung cancer & Lung & 156 & 171 \\
GSE25066 & Breast cancer & Breast & 486 & 165 \\
GSE36961 & Hypertrophic cardiomyopathy & Heart & 145 & 178 \\
GSE39582 & Colorectal cancer & Colon & 560 & 171 \\
GSE42568 & Breast cancer & Breast & 121 & 168 \\
GSE61635 & Systemic lupus erythematosus & Whole blood & 129 & 143 \\
GSE6764 & Hepatocellular carcinoma & Liver & 72 & 178 \\
\bottomrule
\end{tabular}
\end{table*}

\begin{table*}[t]
\centering
\caption{Partitioning algorithms evaluated, grouped by methodological family.}
\label{tab:algos}
\small
\begin{tabular}{ll}
\toprule
Family & Algorithms \\
\midrule
A. Hierarchical co-expression & WGCNA-style (average linkage, dynamic-style cutting) \\
B. Modularity optimisation & Leiden ($\gamma=0.1,0.5,1.0,2.0,5.0,10.0$); Louvain ($\gamma=0.5,1.0,2.0$); \\
 & CPM ($\gamma=0.05,0.10,0.20$); Multilevel; Fast Greedy; Leading Eigenvector; Spin Glass \\
C. Information-theoretic / spectral & InfoMap; Walktrap; spectral clustering \\
D. Generative / label propagation & Label propagation \\
E. Overlapping communities & Clique percolation \\
F. Graph representation learning & Self-supervised GAT; GraphSAGE (edge-prediction objective) \\
G. Classical clustering & $K$-means; Gaussian mixture; Ward; DBSCAN; HDBSCAN; OPTICS; \\
 & Mean Shift; Affinity Propagation; NMF \\
\midrule
\multicolumn{2}{l}{\emph{Total: 31 algorithms from 7 families}} \\
\bottomrule
\end{tabular}
\end{table*}

%% ============================================================
\section{Discussion}
%% ============================================================

\subsection{Principal findings}

We conducted a paired, multi-cohort benchmark of 31 partitioning algorithms on identical
pathway-derived graphs, evaluated across six independent dimensions with synthetic
ground-truth controls. Four conclusions follow. First, multi-resolution community detection
yields functionally more coherent partitions than single-scale hierarchical cutting, without
sacrificing modularity (Section~\ref{sec:main}). Second, synthetic benchmarks with known
nested structure show that the advantage is attributable to an actual resolution capability:
multi-resolution methods recover fine-grained sub-groups that hierarchical cutting merges.
Third, the resolution parameter must be chosen against an external functional criterion;
module count is not a valid proxy for quality. Fourth, and most consequentially for practice,
modularity does not exceed degree-preserving null models on these networks
(Section~\ref{sec:null}), and self-supervised embeddings only partially reproduce the
inferred communities.

\subsection{Implications for practice}

Pathway-centric module analysis is common in transcriptomic studies, and the choice of
partitioning method is rarely justified. Our results support three practical
recommendations. (i) \emph{Report functional coherence against a pathway-internal null.}
Enrichment $p$-values are biased by module size and, within a functionally coherent pathway,
are inflated for any partition; a size-matched, pathway-internal random baseline removes
both biases. (ii) \emph{Do not use module count or modularity alone as evidence of quality.}
Modularity failed to exceed degree-preserving null models, and module count is trivially
controlled by the resolution parameter. (iii) \emph{Select resolution against an external
criterion} --- functional coherence, stability across resamples, or downstream predictive
performance --- rather than by convention.

\subsection{Limitations}

The analysis is restricted to seven GO biological-process pathways and eleven cohorts;
although this exceeds typical practice, it remains modest by the standards of large-scale
benchmarking consortia. Density-based algorithms (DBSCAN, OPTICS, Mean Shift) were evaluated
after spectral embedding to mitigate the curse of dimensionality, which may understate their
performance on raw similarity profiles. Four high-complexity algorithms (edge betweenness,
link communities, optimal modularity and full multi-scale Markov Stability) could not be
evaluated reliably at the pathway sizes considered and were excluded; this is a limitation of
feasibility rather than a negative finding about those methods. Finally, the downstream
evaluation is limited to disease-status classification, which is an easy task in these
cohorts; survival and drug-target endpoints would be more discriminative.

\subsection{Future work}

Three directions follow naturally. First, extending the benchmark to single-cell and
multi-omic data would test whether the conclusions hold when co-expression is estimated from
sparser, more heterogeneous measurements. Second, overlapping community detection ---
included here only in a limited form --- deserves systematic evaluation, since pleiotropy
implies that hard partitions may be an inappropriate model for gene function. Third,
replacing pathway-referenced benchmarks with experimentally derived complexes (e.g. CORUM)
would provide a more direct biological ground truth.

%% ============================================================
\section{Conclusion}
%% ============================================================

On identical pathway-derived co-expression graphs, multi-resolution community detection
produced functionally more coherent modules than hierarchical co-expression clustering,
recovered nested functional sub-structure in synthetic benchmarks with known ground truth,
and yielded better downstream clinical prediction --- while modularity failed to exceed
degree-preserving null models and self-supervised embeddings only partially reproduced the
inferred communities. We therefore recommend that pathway-centric module analyses report
functional coherence against pathway-internal null expectations and select resolution
parameters against external criteria, rather than relying on module counts or modularity.

%% ============================================================
\section*{Acknowledgements}
%% ============================================================

The authors thank [to be completed] for discussions. Computational resources were provided
by [to be completed].

\section*{Funding}
[To be completed.]

\section*{Data availability}
All expression cohorts are publicly available from GEO under the accessions in
Table~\ref{tab:datasets}. Processed matrices, complete evaluation tables (2170
main-benchmark, 528 synthetic-benchmark and 317 downstream records), statistical outputs and
analysis scripts are organised as a reproducible pipeline; code is available at [repository].

\bibliographystyle{unsrtnat}
\bibliography{refs}

\begin{figure}[t]
\centering
\includegraphics[width=\columnwidth]{../output/figures_v3/Fig8_network_structure.png}
\caption{Network structure of an apoptotic-process subnetwork (GSE137342, 110
highest-variance genes) partitioned on an identical topological-overlap input. Node colour
denotes module membership; node size scales with degree. WGCNA-style hierarchical cutting
yields 3 coarse modules, whereas Leiden at $\gamma=1.0$ and $\gamma=2.0$ resolve 4 and 23
modules respectively. Edges are the top 8\% of topological-overlap weights.}
\label{fig:network}
\end{figure}

\begin{figure}[t]
\centering
\includegraphics[width=\columnwidth]{../output/figures_v3/Fig9_umap_embedding.png}
\caption{Self-supervised embedding of the apoptotic subnetwork, projected by UMAP. Nodes
are coloured by partition membership for (a) WGCNA-style hierarchical cutting,
(b) Leiden $\gamma=1.0$ and (c) Leiden $\gamma=2.0$. (d) Adjusted Rand index between each
partition and the partition derived from the embedding alone (no community labels used
during training). The embedding reproduces the Leiden $\gamma=1.0$ partition several-fold
better than the hierarchical one.}
\label{fig:umap}
\end{figure}

\end{document}
